\chapter{Results and Discussion}
\label{ch:results}


\section{Fouling profiles of the as-received membranes}
\label{sec:res-fouling}
This section reports the condition of the two EoL membrane elements before any treatment, as determined by the Membrane Works autopsies conducted \citep{MembraneWorks2025Desal,MembraneWorks2026Mine}. … All three elements were lightly fouled. 4238A, the first-pass desal element, was in good condition and still performed to specification. 4238C, the second-pass desal element, mainly showed damage to the active layer with loss of separation performance. 4297B had light particle and biological fouling, but its polyamide layer was mostly intact.

\subsection{Desal-site first-pass element (4238A)}
\label{sec:red-fouling-4238A}
\paragraph{Desposits.}
4238A was in good condition. A very light brown fouling layer was visible when membrane was wiped, and the membrane had a yellow-to-pink colour that did not wipe off. The brine seal and permeate tube were clean, there was no signs of delamination or backpressure damage, and permeate side was clean \citep{MembraneWorks2025Desal}.

SEM-EDS found iron mixed with aluminosilicate (clay) particles, up to about 13 at.\% Fe, 12 at.\% Si and 11 at.\% Al. Regions without deposits showed mainly C, O and S from the polysulfone support. The FTIR spectra of used 4238A and new 4238B were very similar.

\paragraph{Active-layer condition.}
4238A still met its specification as an element: normalised flux was \qty{26.6}{LMH} against \qty{22.1}{LMH} expected, and salt rejection was \qty{99.65}{\percent} against \qty{99.68}{\percent} expected \citep{MembraneWorks2025Desal}. In the dye test it showed only minor pinholing along the lines of brine-spacer contact, similar to the new elements, with \qty{91}{\percent} rejection in the dye cell. Its contact angle (\qty{25.1}{\degree} against \qty{44.2}{\degree} for new 4238B) and zeta potential at pH~7 (\qty{-30}{mV} against \qty{-8}{mV}) were both different from the new element. Membrane Works attributed the lower contact angle to adsorbed organic or biological material \citep{MembraneWorks2025Desal}.

\subsection{Desal-site second-pass element (4238C)}
\label{sec:res-fouling-4238C}

\paragraph{Deposits.} 
The 4238C membrane was visually clean, with a slight yellow colour. Wiping and scraping the surface removed a small amount of brown deposit. The end caps, brine seal and permeate tube were clean, and there was no damage on the permeate side \citep{MembraneWorks2025Desal}. 

SEM-EDS found three types of deposits:
\begin{itemize}
  \item Iron oxide was found in patches. Some regions showed up to 30.9\% Fe with small amounts of C, S, Al and Si. 
  \item Biological materials. Regions contain up to 14.3\% of nitrogen content. These also contained adsorbed inorganic salts.
  \item Sodium chloride crystals (45.4\% Na, 38.7\% Cl).
\end{itemize}
Areas without any deposits comprised mainly of C, O and S, which comes from the polysulfone support layer. 

The ATR-FTIR spectra of 4238C, and the brand new 4238D were very similar, and no significant chemical difference between the new and used membranes were detected \citep{MembraneWorks2025Desal}. 

\paragraph{Active layer condition}
The main problem with 4238C was damage to the active layer, not fouling. SEM showed holes in the active layer across the surface, and energy dispersive spectroscopy (EDS) couldn't identify a cause. Similar but smaller holes were also observed on the unused 4238D.

The dye test also supports this. Of the four desalination site membranes tested, 4238C let the most dye through, with dense spots and streaks of dye visible across its permeate side (Figures~\ref{fig:dye_test_full} and~\ref{fig:dye_test_magnified}). Its permeate conductivity during the dye test was \qty{113}{\micro\siemens\per\centi\metre}, compared with around \qtyrange{20}{30}{\micro\siemens\per\centi\metre} for the others. By comparison, the used first-pass element 4238A showed low dye penetration, similar to both new elements (\qty{91}{\percent} rejection in the dye cell).More dye also stuck to the feed side of 4238C, which can indicate fouling or surface damage (Figure~\ref{fig:dye_test_full}).

% =====================================================================
% Dye testing figures — adapted from the Membrane Works desalination
% autopsy report. Requires in preamble:
%   \usepackage{graphicx}
%   \usepackage{array}
%   \usepackage{booktabs}
% Images: figures/dye_test/
% =====================================================================

% ---------------------------------------------------------------------
% Figure 1: Whole coupons, feed vs permeate (2nd-pass elements)
% ---------------------------------------------------------------------
\begin{figure}[htbp]
  \centering
  \setlength{\tabcolsep}{3pt}
  \renewcommand{\arraystretch}{1.1}
  \begin{tabular}{>{\centering\arraybackslash}m{0.16\linewidth}
                  >{\centering\arraybackslash}m{0.38\linewidth}
                  >{\centering\arraybackslash}m{0.38\linewidth}}
    \toprule
    \textbf{Membrane} & \textbf{Feed side} & \textbf{Permeate side} \\
    \midrule
    Used 2\textsuperscript{nd} pass \newline (4238C) &
      \includegraphics[width=\linewidth]{figures/dye_test/full_4238C_feed} &
      \includegraphics[width=\linewidth]{figures/dye_test/full_4238C_permeate} \\[2pt]
    New 2\textsuperscript{nd} pass \newline (4238D) &
      \includegraphics[width=\linewidth]{figures/dye_test/full_4238D_feed} &
      \includegraphics[width=\linewidth]{figures/dye_test/full_4238D_permeate} \\
    \bottomrule
  \end{tabular}
  \caption{Dye integrity test of desalination site membrane coupons, showing the feed side (left) and permeate side (right). Dye on the permeate side indicates passage through defects in the polyamide layer. Adapted from \cite{MembraneWorks2025Desal}.}
  \label{fig:dye_test_full}
\end{figure}

% ---------------------------------------------------------------------
% Figure 2: Magnified images, all four elements
% ---------------------------------------------------------------------
\newlength{\dyeh}\setlength{\dyeh}{3.6cm}  % one knob for row height
\begin{figure}[htbp]
  \centering
  \setlength{\tabcolsep}{2pt}
  \renewcommand{\arraystretch}{1.0}
  \begin{tabular}{>{\centering\arraybackslash}m{0.13\linewidth}
                  >{\centering\arraybackslash}m{0.18\linewidth}
                  >{\centering\arraybackslash}m{0.18\linewidth}
                  >{\centering\arraybackslash}m{0.40\linewidth}}
    \toprule
    \textbf{Membrane} & \textbf{Feed side} & \textbf{Permeate side} & \textbf{Damaged areas} \\
    \midrule
    Used 1\textsuperscript{st} pass \newline (4238A) &
      \includegraphics[height=\dyeh]{figures/dye_test/mag_4238A_feed} &
      \includegraphics[height=\dyeh]{figures/dye_test/mag_4238A_permeate} &
      \includegraphics[height=\dyeh]{figures/dye_test/mag_4238A_damage1} \\[2pt]
    New 1\textsuperscript{st} pass \newline (4238B) &
      \includegraphics[height=\dyeh]{figures/dye_test/mag_4238B_feed} &
      \includegraphics[height=\dyeh]{figures/dye_test/mag_4238B_permeate} &
      \includegraphics[height=\dyeh]{figures/dye_test/mag_4238B_damage1}\hspace{2pt}%
      \includegraphics[height=\dyeh]{figures/dye_test/mag_4238B_damage2} \\[2pt]
    Used 2\textsuperscript{nd} pass \newline (4238C) &
      \includegraphics[height=\dyeh]{figures/dye_test/mag_4238C_feed} &
      \includegraphics[height=\dyeh]{figures/dye_test/mag_4238C_permeate} &
      \includegraphics[height=\dyeh]{figures/dye_test/mag_4238C_damage1}\hspace{2pt}%
      \includegraphics[height=\dyeh]{figures/dye_test/mag_4238C_damage2} \\[2pt]
    New 2\textsuperscript{nd} pass \newline (4238D) &
      \includegraphics[height=\dyeh]{figures/dye_test/mag_4238D_feed} &
      \includegraphics[height=\dyeh]{figures/dye_test/mag_4238D_permeate} &
      \includegraphics[height=\dyeh]{figures/dye_test/mag_4238D_damage1}\hspace{2pt}%
      \includegraphics[height=\dyeh]{figures/dye_test/mag_4238D_damage2} \\
    \bottomrule
  \end{tabular}
  \caption{Magnified views of the dye integrity test for all four elements, with close-ups of localised damage. Adapted from \cite{MembraneWorks2025Desal}.}
  \label{fig:dye_test_magnified}
\end{figure}


Standard performance testing by \cite{MembraneWorks2025Desal} on a single 8" element showed the effect of the membrane damage. The test used a feed salt concentration of \qty{1500}{mg/L} at a pressure of \qty{10}{bar}. The 4238C EoL membrane had a flux of \qty{71.9}{LMH}, about twice the expected \qty{37.3}{LMH}, and a salt rejection of \qty{97.3}{\percent}, below the expected \qty{99.46}{\percent}. High flux combined with low rejection indicates that the active layer may be chemically or physically damaged.

\flag{EXPLAIN}
The surface properties of 4238C had also changed compared to the new 4238D, even though FTIR showed no chemical difference. The isoelectric point increased from pH~3.1 to 3.6; The zeta potential at pH~7 became less negative, from \qty{-42}{mV} to \qty{-16}{mV}; the water contact angle decreased from \qty{58.5}{\degree} to \qty{34.5}{\degree}, meaning the surface became more hydrophilic.

The lower contact angle could be attributed to organic or biological material adsorbed onto the surface, and the loss of negative charge is similar to what has been reported for membranes cleaned repeatedly with caustic \cite{MembraneWorks2025Desal}. The pH~10 feed to the second pass may have had a similar effect. A less negative surface repels borate ions less strongly, which is consistent with the falling boron rejection reported by the plant.


\subsection{Mine-site element (4297B)}
\paragraph{Deposits.}
The mine-site feed water is high in calcium and sulfates. Using water chemistry modelling, calcium sulphate (gypsum) and silica were the main scaling risk in the feed for this membrane \citep{MembraneWorks2026Mine}. The 4297B membrane had a light brown fouling layer across the surface. Both end caps had spots of brown fouling but no debris, and brine spacer was clean. 

Although there was significant visible fouling on the membrane, fouling analysis by extraction was conducted and showed that the amount of extractable foulant was low. Around \qty{237.5}{mg.m^{-2}} with similar composition across the element: 

\begin{itemize}
  \item calcium: \qtyrange{91}{92}{mg.m^{-2}};
  \item carbon: \qtyrange{69}{75}{mg.m^{-2}};
  \item magnesium: \qtyrange{30}{31}{mg.m^{-2}};
  \item sulfur, likely as sulfate: \qtyrange{26}{27}{mg.m^{-2}};
  \item silicon: \qtyrange{10}{11}{mg.m^{-2}};
  \item iron: \qty{1}{mg.m^{-2}} or less.
\end{itemize}

SEM-EDS showed particles scattered across the surface of the membrane. Analysis detected iron oxide (37.3 \% Fe and 50.7\% O, with small amounts of Al, S and Si (2.8, 2.38, 1.81\% respectively)) in some regions, as well as silica (29.6\% Si and 65.8\% O) in others. A low level of nitrogen (5.4 \%) was also seen across the surface, indicating a small amount of biofouling. 

While Iron was visible in the SEM-EDS, almost non was recovered by extraction, suggesting that iron oxide was not easily dissolved by acid or caustic. This is discussed further in Section~\ref{sec:res-cleaning}

4297B also showed more signs of biological fouling than 4238C. Astra Blue staining showed brown fouling under the brine spacer, attributed to extracellular polymeric substance. The bacterial growth count was \qtyrange{5800}{7800}{CFU/mL} classifying as moderate growth \citep{MembraneWorks2026Mine}.

\paragraph{Active-layer condition.}
A swab of the permeate tube showed slight fouling and the back of the membrane was stained brown. Under the optical microscope, the foulant appeared to have penetrated into the support layers. SEM also showed indentations in the surface where the scale had previously been. 

The polyamide layer of 4297B showed minor chlorine damage, with a negative Fujiwara test (which determines polyamide RO membrane exposure to oxidising halogens). XPS detected \qty{0.2}{\percent} organic chlorine, indicating some low level chlorine damage is present. Incorrect sodium metabisulfite (SMBS) dosing can react with iron on the membrane and catalyse SMBS degradation and cause oxidative damage \citep{MembraneWorks2025Desal}.

Performance of the 4297B membrane was not tested by Membrane Works, and the reported drop to around \qty{85}{\percent} rejection was measured for the whole RO train. The autopsy report concluded that the fouling and damage did not fully explain this drop. 




\subsection{Comparison of EoL membranes}
Both 4238A and C came from the same plant and served the same \qty{4.5}{years}, but 4238A sat in the first pass at neutral pH and showed light fouling with no loss of performance. The damage to 4238C therefore cannot be attributed to the age of the plant or the plant's feed water. It is more likely linked to its second-pass duty, including the pH~10 feed.
\flag{ AI GENERATED SUMMARY TABLE-- CHECK}
\begin{table}[htbp]
  \centering
  \small
  \caption{Fouling profiles of the as-received elements, compiled from
  \cite{MembraneWorks2025Desal,MembraneWorks2026Mine}.}
  \label{tab:fouling-profiles}
  \begin{tabular}{p{3.2cm}p{5.4cm}p{5.4cm}}
    \toprule
    & 4238C (desal) & 4297B (mine) \\
    \midrule
    Service & \qty{4.5}{y}; 2nd pass, feed pH~$\approx$10
            & \qty{2.75}{y}; single stage, high-sulfate feed \\
    Visual deposit & Slight yellowness; slight brown deposit on wiping
                   & Light brown layer; back-side staining \\
    Dominant deposits (SEM-EDS) & Fe/Al oxide; N-rich biological particles;
                                  NaCl crystals
                                & Fe oxide; silica (sand); trace biofouling \\
    Extractable loading & --- & \qtyrange{237}{238}{mg.m^{-2}} (Ca, C, Mg, S) \\
    Biological activity & N-rich particles only
                        & Astra Blue positive; 5,800~CFU\,mL$^{-1}$ \\
    Physical damage & Holes (SEM); most dye passage of set
                    & Indentations from removed scale \\
    Chemical change & FTIR unchanged; IEP 3.1\,$\to$\,3.6;
                      $\zeta_{\mathrm{pH\,7}}$ \qty{-42}{}\,$\to$\,\qty{-16}{mV};
                      contact angle \qty{58.5}{}\,$\to$\,\qty{34.5}{\degree}
                    & Fujiwara negative; \qty{0.2}{\percent} organic Cl;
                      N retained (XPS) \\
    As-received performance & \qty{71.9}{LMH}, \qty{97.3}{\percent}
                              (expected \qty{37.3}{LMH}, \qty{99.46}{\percent})
                            & --- \\
    \bottomrule
  \end{tabular}
\end{table}

Both membranes were lightly fouled, and both had significant iron oxide and biological material on the surface. The main difference between the two is that 4238C had damage to the active layer, where holes were seen in SEM and had a higher dye passage. This membrane had double the expected flux and lower rejection. 4297B on the other hand had much more deposits, such as calcium, magnesium and sulfate in the extracts, silica, iron oxide that was difficult to remove by conventional cleaning methods, and a biological film. The polyamide layer showed minor chlorine damage and indentation from scaling. 

Different tests were run on each membrane depending on the reported issue and diagnostic methods. Only visual inspection and SEM-EDS was done on both. 

\subsection{Implications for regeneration}
\flag{link points to section \ref{sec:fouling-ip}}

The two membrane are expected to respond differently to a new IP layer. 

For 4238C, the surface is mostly clear of foulants and deposits, but the new polyamide layer has to cover existing holes in the active layer. The change in surface charge and hydrophilicity can also affect how the surface takes up the aqueous monomer as compared with traditional IP methods.

For 4297B, the deposits are more likely to affect IP. Iron oxide and silica particles here did not appear in the extraction sample from caustic soda and acid, but do appear on SEM-EDS analysis, suggesting the compounds arn't easily broken down by acid and base cleans. The biological film can also change the surface that the new layer forms on. As a result, it is likely that despite extensive cleaning of membranes, these foulants still remain right up till the IP process of uplifting.

%-------------------––––Next Section-----------------%
\section{Performance of the as-recieved membranes}
\label{sec:res-baseline}
Uncleaned coupons without IP were tested on the crossflow rig to give a baseline for each membrane (Table~\ref{tab:baseline}). Later results are compared against these values.

  \begin{table}[htbp]
  \centering
  \caption{Performance of the as-received membranes (no clean, no IP). Mean $\pm$ SD across coupons.}
  \label{tab:baseline}
    \begin{tabular}{llcccc}
    \toprule
    Campaign & Membrane & Flux (\unit{LMH}) & PWP (\unit{LMH.bar^{-1}}) & NaCl rejection (\%) & $n$ \\
    \midrule
    1 & 4238A (desal, 1st pass) & $1.23 \pm 0.27$ & $0.082 \pm 0.018$ & $78.1 \pm 2.8$ & 2 \\
    1 & 4238C (desal, 2nd pass) & $0.56 \pm 0.08$ & $0.038 \pm 0.006$ & not measured & 2 \\
    2 & 4238C (desal, 2nd pass) & $3.21 \pm 0.32$ & $0.208 \pm 0.011$ & $76.6 \pm 3.1$ & 4 \\
    2 & 4297B (mine)            & $5.22 \pm 1.00$ & $0.330 \pm 0.062$ & $81.9 \pm 8.6$ & 4 \\
    \bottomrule
  \end{tabular}
  \par\smallskip
  \footnotesize $n$ = number of coupons. Flux and PWP at \qtyrange{14}{16}{bar}; rejection with $\sim$\qty{2000}{mg.L^{-1}} NaCl (\qtyrange{3800}{4000}{\micro S.cm^{-1}}).
  \end{table}

In Campaign~2, PWP was higher for the mine-site element than for 4238C (0.330 against \qty{0.208}{LMH.bar^{-1}}; Welch's $t$-test, $p = 0.028$), while rejection did not differ significantly (81.9 against \qty{76.6}{\percent}, $p = 0.31$). The as-received PWP of 4238C differed by a factor of about five between campaigns (0.038 against \qty{0.208}{LMH.bar^{-1}}), although both sets of coupons came from the same element. Within each campaign the spread between coupons was small (RSD 15\,\% and 5\,\%). The difference therefore lies between regions of the membrane sheet sampled in each campaign, not in the measurement itself. This is discussed below.

\subsection{Comparison with element test}
4238A testing becomes a useful check, because the EoL element still met the specification as an element (\qty{99.65}{\percent}) rejection. During coupon testing however, the \flag{What was tested PWP?} PWP was only measured to be \qty{0.082}{LMH.bar^{-1}} and a rejection of \qty{78.1}{\percent}, which is 21 percentage points below the element testing results. This gap is similar for the damaged 4238C (about 21~pp). This means taht the difference between coupon and element performance cannot simply be explained by the damage to 4238C. 

For 4238C, the values are much lower than those measured by Membrane Works on the full elements. Their test gave a flux of \qty{71.9}{LMH} at about \qty{10}{bar}, or roughly \qty{7}{LMH.bar^{-1}}, with a rejection of \qty{97.3}{\percent} \citep{MembraneWorks2025Desal}. The coupon permeability is about 35 times lower, and the rejection about 21 percentage points lower. The permeability measured in this experiment is about 21 times lower than the rated value for a new ESPA2-MAX element (about \qty{4.4}{LMH.bar^{-1}}). 

For 4297B membranes, no testing was done on the RO element, and the reported drop to around \qty{85}{\percent} was measured for the whole RO train \citep{MembraneWorks2026Mine}, so there are no element values to compare with. 

\subsection{Why coupon results differ}
Testing using the 8-inch RO module from Membrane Works revealed a much higher permeability and rejection, at \qty{71.9}{LMH} flux, 97.3\% rejection at 10 bar and \qty{1500}{mg/L}. That experiment used a second-pass ESPA2-MAX element which was reported to be underperforming for desalination purposes. The flux measured on site of that element was double the expected \qty{37.3}{LMH}, ane below the expected 99.46\%, likely due to pinhole damage in the active layer. Laboratory flat-sheet dye-cell test also gave a rejection of only 61\%, compared with 91-92\% for other coupons tested.
4238A however shows that this explanation is incomplete, were performance as an element yielded \qty{99.65}{\percent} on the rig, essentially the same as the damaged 4238C (\qty{76.6}{\percent}). Coupon PWP was also 15-35 times lower than the corresponding elements for all three membranes. Two things could have contributed to this difference. First is the idea of apparent rejection, where at constant salt permeability, the observed rejection falls as water flux falls. Coupon fluxes of \qtyrange{0.5}{6}{LMH} are an order of manitude below the element fluxes, so much lower observed rejection is expected (given no change to salt permeability). It is also possible that salt could have leaked through via local defects or O-ring seal defects. A small coupon area of \qty{0.0139}{m^2} makes this leak more apparent. The small coupon size also means local defects in coupon is amplified greatly, rather than averaged out like it would for an element testing. 

Because this offset applied to every condition tested on the rig, comparisons between the conditions are valid, but these values should not be compared to the element or datasheet values as there is clearly a big disparity between them. 

%-------------------––––Next Section-----------------%
\section{Cleaning outcomes}
\label{sec:res-cleaning}
This section reports what the cleaning steps removed from each membrane. The uncleaned EoL were compared with the extensive clean, as described in Section~\ref{sec:cleaning}. The effect of cleaning on regeneration is covered in Section~\ref{sec:fouling-ip}.

\subsection{Fouling profiles and expected cleaning response}
\flag{Repeats alot of what was said in 4.1. Cross refernce to section \ref{sec:res-fouling}}\\
The autopsies give a clear idea of what cleaning could remove.

For 4238C, there was little deposit across the membrane. The membrane was visually clean, and SEM-EDS found scattered iron and aluminium oxide, biological particles and salt crystals. Cleaning was therefore not expected to produce a large visible change. The main issues found with 4238C were the holes in the active layer and changes in surface chemistry and charge, which cleaning cannot repair.

For 4297B, the deposit comprised different elements and compounds, all expected to react to cleaning differently.
\paragraph{Calcium, magnesium and sulfate.} These made up most of the extractable foulant and were recovered by acid and caustic extraction in the fouling extraction test conducted by Membrane Works. They were expected to be removed in a similar fashion by acid clean. EDTA, which binds metal ions such as calcium and magnesium, was expected to help as its chelating ability increases at alkaline pH where more of its carboxylic groups are deprotonated \citep{elimelech2005chemical}.
\paragraph{Organic and biological materials.} The extractable carbon and biofilm seen with Astra Blue staining were expected to be removed mainly by the caustic clean. However, where organic foulants are cross-linked by calcium, NaOH alone has been shown to give poor flux recovery, while EDTA was considerably more effective by breaking the calcium bridges \cite{elimelech2005chemical}.
\paragraph{Iron oxide.} \flag{Need to reevaluate, cannot assume no iron oxide cuz no colour change} Almost no iron was recovered by acid and caustic extraction, even though iron oxide was clearly visible in SEM-EDS. The iron oxide was therefore not expected to be fully removed by HCl and NaOH cleaning.
\paragraph{Silica particles.} These were identified by Membrane Works as sand rather than scale, and were not expected to dissolve in either solution.\\

Additionally, Membrane Works conducted a cleaning study supporting a lot of the points above. A high pH biological and clay cleaner showed some improvements to the fouling profile with the advanced iron cleaning and specialty silica cleaning showing similar results. Increasing the advanced iron cleaner from \qty{2}{\percent} to \qty{4}{\percent} removed more of the fouling, but some remained. Since specialist iron cleaners did not fully remove the iron, the HCl and NaOH cleans here were expected to leave iron oxide on 4297B.

% ---------------------------------------------------------------------
\subsection{Expected effect of extended cleaning on membrane performance}
Cleaning chemicals can remove foulants in membrane but can also alter the polyamide layer itself. Change in PWP or salt rejection can be cause by the cleaning itself, rather than simply foulant removal. \cite{tu2015chemical} evaluated this effect by soaking virgin ESPA2 low pressure RO membranes in cleaning solutions for \qty{25}{\hour} at \qty{30}{\celsius}. This is expected to represent the cumulative cleaning over about three years of plant operation. As no foulants were present in this study, any performance could be attributed to the chemical interacting with the membrane.

\paragraph{Water permeability.} Solution pH was the main driver of the permeability change. NaOH at pH~11 increased PWP by about \qty{5}{\percent}, while citric acid at pH~3 decreased it by more than \qty{10}{\percent}\citep{tu2015chemical}. The caustic increase was attributed to repulsion between  

\paragraph{Salt rejection.} 

\paragraph{Implications for this study.} 



\subsection{Visual observations}
\paragraph{Desal-site membrane 4238C. } 
Cleaning caused no visible change to 4238C. Before cleaning, the polyamide (feed side) was a shiny, light yellow colour, and the polyester (permeate side) was a dull, matte white. After both the acid and base cleans, the coupon looked exactly the same. The acid and base solutions also showed no colour change after use. This suggests that little foulant was dissolved into the solutions, which is expected according to the autopsy as only light deposits were seen.

\paragraph{Mine-site membrane 4297B. }
The colour of 4297B did not change noticeably after cleaning. \flag{The acid and base solutions also showed no colour change, which indicates that the iron oxide was not dissolved. This agrees with the autopsy, where acid and caustic extraction recovered almost no iron.}

The main observation for 4297B was the physical damage at the edges of the coupons (Figure~\flag{1}). In some areas, the polyamide layer had started to peel away from the polysulfone support and could be lifted off with very little force. This was first seen immediately after the base clean, and the damage appeared in repeated experiments too. No delamination was seen on 4238C under the same cleaning conditions, and the autopsy also found no delamination on the EoL 4297B elements. 

The base clean reached pH of up to 11.5, which is slightly outside the manufacturer's limit for continous operation (2-11 pH) but within the short-term (30 min) cleaning range (1-13) \citep{MembraneWorks2026Mine}. The potential damage of membranes from intensive clean outside of recommended pH ranges is explored in depth in Section~\ref{sec:res-cleaning-chem}. One possible explanation is that the bond between the polyamide and support was already weakened during service. The autopsy found that foulants appeared to have penetrated into the support layers of 4297B, and indentations on the surface from past scale could contribute to this. The damage was also concentration along the edges of the coupon, suggesting the cut may have exposed the gap between the layers to the cleaning solution, allowing it to penetrate through. 

Delaminated areas have no polyamide layer present, so they would expect to have UF/MF like performance (high permeability and low rejection) if included in the test. Instead, the damaged edges of the coupon were trimmed and discarded.

\flag{add image here}

\subsection{Performance after cleaning}
\label{sec:res-cleaning-perf}
%PWP and rejection after extensive clean, compared with Section~\ref{sec-baseline} for both membranes

Table~\ref{tab:clean-compare} shows the flux, PWP and NaCl rejection of the cleaned coupons. As-recieved values are included for reference and comparison.
\begin{table}[htbp]
  \centering
  \caption{As-received versus extensively cleaned coupons (no IP). Mean $\pm$ SD across coupons; change with Welch's $t$-test $p$-value.}
  \label{tab:clean-compare}
  \small
  \begin{tabular}{lllccc c}
    \toprule
    Campaign & Membrane & Condition & Flux (\unit{LMH}) & PWP (\unit{LMH.bar^{-1}}) & Rejection (\%) & $n$ \\
    \midrule
    1 & 4238A & As received & $1.23 \pm 0.27$ & $0.082 \pm 0.018$ & $78.1 \pm 2.8$ & 2 \\
      &       & Cleaned     & $2.81 \pm 1.02$ & $0.198 \pm 0.072$ & $86.9 \pm 9.2$ & 2 \\
      &       & Change ($p$) & & $+142\%$ (0.25) & $+8.8$\,pp (0.39) & \\
    \midrule
    1 & 4238C & As received & $0.56 \pm 0.08$ & $0.038 \pm 0.006$ & not measured & 2 \\
      &       & Cleaned     & $13.10 \pm 0.35$ & $0.922 \pm 0.025$ & $96.2 \pm 1.1$ & 2 \\
      &       & Change ($p$) & & $\times 24$ (0.009) & --- & \\
    \midrule
    2 & 4238C & As received & $3.21 \pm 0.32$ & $0.208 \pm 0.011$ & $76.6 \pm 3.1$ & 4 \\
      &       & Cleaned, all & $4.55 \pm 4.42$ & $0.295 \pm 0.290$ & $79.7 \pm 9.7$ & 7 \\
      &       & Change ($p$) & & $+42\%$ (0.46) & $+3.1$\,pp (0.46) & \\
      &       & Cleaned, excl.\ C5--C6 & $2.04 \pm 1.16$ & $0.130 \pm 0.077$ & $84.3 \pm 6.9$ & 5 \\
      &       & Change ($p$) & & $-37\%$ (0.086) & $+7.7$\,pp (0.069) & \\
    \midrule
    2 & 4297B & As received & $5.22 \pm 1.00$ & $0.330 \pm 0.062$ & $81.9 \pm 8.6$ & 4 \\
      &       & Cleaned     & $5.60 \pm 1.41$ & $0.359 \pm 0.073$ & $80.5 \pm 11.6$ & 6 \\
      &       & Change ($p$) & & $+9\%$ (0.52) & $-1.4$\,pp (0.83) & \\
    \bottomrule
  \end{tabular}
  \par\smallskip
  \footnotesize Flux and PWP at \qtyrange{14}{16}{bar}; rejection with $\sim$\qty{2000}{mg.L^{-1}} NaCl. Campaign~1 cleaned conditions were tested in a single run (two coupons, one rejection reading each).
\end{table}

  For the mine-site membrane (4297B), the six cleaned coupons had fairly consistent flux readings. The flux ranged from \qty{3.79}{LMH} to \qty{7.22}{LMH}, giving a mean of 5.60 $\pm$ \qty{1.41}{LMH}. Rejection varied significantly from \qty{64.2}{\percent} to \qty{91.2}{\percent}, mainly because two coupons tested yielded a much lower rejection than others. 

  The desalination plant showed much more variation between the tests. Across seven coupons, the mean flux was 4.55 $\pm$ \qty{4.42}{LMH}, with a relative standard deviation of about \qty{97}{\percent}. This is a huge spread of readings, but does not come from pure measurement error. Taking repeated readings on the same coupon differed by at most \qty{0.34}{LMH}, while coupons had a mean flux of \qty{0.67}{LMH} to \qty{11.76}{LMH}. Most of the spread came from coupon sample 5 and 6 (C5 and C6 from hereon), where they were measured to have a flux of \qty{9.86}{LMH} and \qty{11.76}{LMH}. These were three to six times higher than the other five coupons, and the rejection of these two coupons were also the lowest measured at \qty{69.3}{\percent} and \qty{67.1}{\percent}.\\

  These two coupons should not be treated as measurment outliers. Rig or test-day effect can also be ruled out, as C7 was tested in the same session at the same pressure, and gave a result close to others (\qty{3.87}{LMH}, \qty{82.7}{\percent}). 4238C also showed no signs of delamination, and permeability readings were taken in ~15 mins interval with no signs of significant change, ruling out compaction damage. The two coupons had both high flux and low rejection, which suggests that feed water leaked through openings in the active layer. 

  The membrane autopsy found similar damage for this membrane. SEM showed holes in active layer, and the dye test showed dense spots and streaks of dye spots spread unevenly across the permeate side as seen in Figure~\ref{fig:dye_test_full} \citep{MembraneWorks2025Desal}. C5 and 6 were cut from the same section of the membrane sheet, so they likely came from regions with a cluster of these pinholes. As the tested membrane coupons are small (\qty{0.0139}{m^2}), these small, local defects can dominate readings in a way that wouldn't happen with full membrane elements. It is highly likely therefore that C5 and 6 are results from a damaged part of the membrane, and do not represent the effect of cleaning. 
  
  Without C5 and 6 coupons, the cleaned 4238C membranes have a mean flux of 2.04 $\pm$ \qty{1.16}{LMH}, a PWP of 0.13 $\pm$ \qty{0.08}{LMH.bar^{-1}} and a rejection of 84.1 $\pm$ \qty{6.5}{\percent} across 5 coupons. The relative standard deviation for flux fell from \qty{97}{\percent} to \qty{57}{\percent}. There is still some spread that remains, but none of the other coupons showed a high-flux low rejection pattern. The values without C5 and 6 will be used to assess the effect of cleaning and to compare against IP.\\

  \flag{}
  Two of the cleaned 4297B coupons gave low rejection (\qty{64.3}{\percent} and \qty{69.3}{\percent}). These were also the two highest-flux coupons (\qty{7.22}{LMH} and \qty{7.13}{LMH}), about 1.5 times the mean of the remaining four (\qty{4.81}{LMH}). This is a similar but much weaker pattern than C5 and C6, whose flux was around five times that of the other 4238C coupons. Unlike 4238C, there is no independent evidence of holes in the 4297B active layer: SEM showed indentations from removed scale rather than holes, and no dye test was performed \citep{MembraneWorks2026Mine}. Higher flux combined with lower rejection is instead consistent with a looser polyamide layer in these regions, which may reflect local variation in the chlorine damage detected by XPS. These coupons were therefore retained. Excluding them, cleaned 4297B gave a flux of $4.81 \pm 0.90$\,\unit{LMH} and a rejection of $87.3 \pm 5.5$\,\unit{\percent}, which is considered in Section~\ref{sec:res-cleaning-compare}.



\begin{table}[htbp]
  \centering
  \caption{Performance of individual cleaned 4238C (desal) coupons.
  Calculated as the average of repeat readings.}
  \label{tab:clean-desal-coupons}
  \begin{tabular}{lcccc}
    \toprule
    Coupon & $\Delta P$ (\unit{bar}) & Flux (\unit{LMH}) & PWP (\unit{LMH.bar^{-1}}) & NaCl rejection (\%) \\
    \midrule
    C1 & 15.9 & 1.91 & 0.120 & 77.5 \\
    C2 & 15.9 & 2.11 & 0.133 & 78.7 \\
    C3 & 15.8 & 1.65 & 0.105 & 92.4 \\
    C4 & 15.8 & 0.67 & 0.042 & 89.3 \\
    \textbf{C5} & 15.3 & \textbf{9.86} & \textbf{0.644} & \textbf{69.3} \\
    \textbf{C6} & 15.3 & \textbf{11.76} & \textbf{0.713} & \textbf{67.1} \\
    C7 & 15.3 & 3.87 & 0.253 & 82.7 \\
    \midrule
    All coupons ($n = 7$)     & & $4.55 \pm 4.42$ & $0.29 \pm 0.28$ & $79.6 \pm 9.4$ \\
    Excl.\ C5--C6 ($n = 5$)   & & $2.04 \pm 1.16$ & $0.13 \pm 0.08$ & $84.1 \pm 6.5$ \\
    \bottomrule
  \end{tabular}
\end{table}

\subsection{Comparison with the as-received membranes}
 \label{sec:res-cleaning-compare}
For 4297B, the extensive clean made little difference. Flux rose by \qty{6.7}{\percent}, PWP by \qty{8.9}{\percent}, and rejection fell by \qty{1.3}{pp}, all within one standard deviation of the coupons. 
\flag{aligns with wollongong paper}

For 4238C, all seven coupons together suggest an increase in flux of \qty{41.8}{\percent}. Excluding C5 and C6 however, the change reverses, where flux falls by \qty{36.4}{\percent}, PWP by \qty{37.3}{\percent}, and rejection rises by \qty{7.6}{pp}. This aligns with autopsy, that showed 4238C did not have alot of foulants to remove (Section~\ref{sec:res-fouling-4238C}). \flag{finish, does not align with autopsy, possibly due to acid tightening or coupon coupon variability. Welch's t-test?}




\subsection{Surface analysis after cleaning}
\label{sec:res-cleaning-surface}
post-clean FTIR/SEM-EDS versus the as-received baseline in Section~\ref{sec:res-fouling}.

% The 2nd pass EPA2-MAX membranes were more typical of what we would expect. The zeta- potential of the new membrane is quite negative and then became less negative in the used membrane. Literature reports that repeated caustic CIP cycles reduces the negative zeta potential (Figure 33) of used membranes in a chemical aging phenomenon. While these membranes have not been repeatedly cleaned, exposure to the high pH 2nd pass could have the same effect1

%1 Phenomenological Interpretations of Membrane Properties, Following Repeated Chemical Cleaning of an End-of-Life Potable Reuse Reverse Osmosis Element Dominated by Inorganic Fouling. Bilal Abada, Jana Safarik, Kenneth P. Ishida, and Shankararaman Chellam*, ACS EST Engg. 2025, 5, 167−178.


%-------------------––––Next Section-----------------%
\section{Regeneration by IP}
\label{sec:res-ip}
%PWP and rejection for no clean and extensive clean + IP compared with same coupons before


\subsection{Performance after IP}
\label{sec:res-ip-perf}
% Scatter plot: EoLm cleaned and regenerated coupons against new commercial RO and literature NF/UF. 

  \begin{table}[htbp]
  \centering
  \caption{Performance of membranes after IP (clean, IP).
  Mean $\pm$ SD.}
  \label{tab:ip-results}
    \begin{tabular}{lcccc}
    \toprule
    Membrane & Flux (\unit{LMH}) & PWP (\unit{LMH.bar^{-1}}) & NaCl rejection (\%) & $n$ \\
    \midrule
    4238C (desal) & $3.21 \pm 0.32$ & $0.21 \pm 0.01$ & $76.5 \pm 3.6$ & 4 / 4 \\
    4297B (mine)  & $5.24 \pm 1.00$ & $0.33 \pm 0.06$ & $81.8 \pm 8.1$ & 4 / 4 \\
    Clean 4238C (desal) & $4.55 \pm 4.42$ & $0.29 \pm 0.28$ & $79.6 \pm 9.4$ & 7 / 7 \\
    Clean 4297B (mine)  & $5.60 \pm 1.41$ & $0.36 \pm 0.07$ & $80.5 \pm 11.5$ & 6 / 6 \\
    IP 4238C (desal) & $0.35 \pm 0.20$ & $0.023 \pm 0.013$ & $76.2 \pm 7.9$ & 4 / 4 \\
    IP 4297B (mine)  & $0.69 \pm 0.19$ & $0.046 \pm 0.012$ & $80.5 \pm 11.2$ & 2 / 2 \\
    \bottomrule
  \end{tabular}
  \par\smallskip
  \footnotesize $n$ = number of cell measurements for Flux \& PWP / Rejection.
  Flux and PWP at \qtyrange{14}{16}{bar}; rejection with $\sim$\qty{2000}{mg.L^{-1}} NaCl
  (\qtyrange{3800}{4000}{\micro S.cm^{-1}}).
  \end{table}

\subsection{Permeability--rejection tradeoff}
\label{sec:res-tradeoff}
% Scatter plot: EoL, cleaned and regenerated coupons against new commercial RO and literature NF/UF.

\subsection{Surface characterisation of regenerated coupons}
\label{sec:res-ip-surface}
% Post IP sem and FTIR



%---------Discussion---------------%
\section{Effect of cleaning on regeneration}
\label{sec:res-cleaning-effect}
Answers: whether the substrate must be cleaned before IP for regeneration to work at all.

\begin{itemize}
  \item no clean versus extensive clean, for both membranes.
  \item an uncleaned substrate should yield a defective film --- high permeability, low rejection, wide replicate scatter.
  \item does the cleaned arm reach commercial RO territory, or only improve on the uncleaned one? 
  \item Distinguish a real effect from replicate scatter; state  limits of what replicate can support.
  \item If the cleaned arm does not deliver the expected permeability drop and rejection gain, interpret  against the five mechanisms in Section~\ref{sec:fouling-ip}
\end{itemize}


\section{Mine site versus Desal site}
\label{sec:res-a-vs-c}
Attributes the differences between the two membranes to their fouling profiles and service history.

\begin{itemize}
  \item Differences in response attributable to feed and operation condition .
  \item Link back to the fouling profiles in Section~\ref{sec:res-fouling}: does fouling type predict regeneration outcome?
\end{itemize}


\section{Discussion synthesis}
\label{sec:res-synthesis}
Reconciles the results with Chapter~\ref{ch:literature_review}, stating what worked, what did not, and which mechanism best explains the data.

\begin{itemize}
  \item Which of the fouling-interference mechanisms in Section~\ref{sec:fouling-ip} the data supports, and which it rules out.
  \item review results against reported regeneration performance (Section~\ref{sec:uplifting}).
  \item Answer research gap from Section~\ref{sec:research-gap} explicitly.
  \item Note anything unexpected
\end{itemize}


\section{Limitations and sources of error}
\label{sec:res-limitations}
States the constraints on how far these results can be generalised, and where the largest uncertainties sit.

\begin{itemize}
  \item Replicate count, coupon cutting, rig variability, measurement uncertainty.
  \item IP variability and cleaning variability between runs.
  \item coupon rather than module scale.
\end{itemize}


\section{Summary tables}
\label{sec:res-tables}
Consolidates the numerical results into a  reference table, raw data deferred to the appendix.

\begin{itemize}
  \item Summary performance table --- rows: each condition; columns: average flux, average PWP (with/without IP), change in permeability, average rejection (with/without IP), change in rejection.
  \item Raw trial data (every trial's volume, time, flux, conductivity) in Appendix~A.
\end{itemize}